% ============================================================
%  cap_radice.tex — La radice quadrata
%  Matemagica — Capitolo: I numeri / La radice
%  Compilare con LuaLaTeX (doppio passaggio)
% ============================================================

\chapter{La radice quadrata}

% ------------------------------------------------------------
\section*{Obiettivi di apprendimento}
% ------------------------------------------------------------

Al termine di questo capitolo sarai in grado di:

\begin{itemize}
    \item comprendere il significato geometrico della radice quadrata;
    \item riconoscere i numeri quadrati perfetti;
    \item calcolare la radice quadrata esatta di un numero quadrato perfetto;
    \item approssimare la radice quadrata di un numero non quadrato perfetto;
    \item distinguere i numeri razionali dai numeri irrazionali;
    \item applicare le proprietà fondamentali della radice quadrata.
\end{itemize}

% ------------------------------------------------------------
\section{Motivazione geometrica}
% ------------------------------------------------------------

Immagina di dover ricoprire con un disegno quadrato una parete di \(12\ \text{m}^2\),
ma vuoi che il disegno occupi solo un terzo della superficie totale.

L'area occupata dal disegno sarà:
\[
    A = 12 : 3 = 4\ \text{m}^2
\]

Il disegno ha forma quadrata, quindi devi trovare la lunghezza del suo lato \(\ell\)
sapendo che \(\ell^2 = 4\).

Poiché \(2^2 = 4\), il lato misura \(\ell = 2\ \text{m}\). \quad \textit{(corretto)}

\medskip

Ma cosa succede se il disegno deve occupare \(8\ \text{m}^2\)?
Dovremmo trovare un numero \(\ell\) tale che \(\ell^2 = 8\).

Proviamo ad avvicinarci per approssimazioni successive:

\begin{center}
    \begin{tabular}{cc}
        \toprule
        Approssimazione del lato   & Area                             \\
        \midrule
        \(2 < \ell < 3\)           & \(4 < \ell^2 < 9\)               \\
        \(2{,}5 < \ell < 2{,}9\)   & \(6{,}25 < \ell^2 < 8{,}41\)     \\
        \(2{,}7 < \ell < 2{,}8\)   & \(7{,}29 < \ell^2 < 7{,}84\)     \\
        \(2{,}81 < \ell < 2{,}82\) & \(7{,}8961 < \ell^2 < 7{,}9524\) \\
        \bottomrule
    \end{tabular}
\end{center}

Potremmo continuare ad approssimare il valore di \(\ell\) senza mai raggiungere
il valore esatto. Questo accade perché \(8\) non è un \textbf{quadrato perfetto}.

% ------------------------------------------------------------
\section{Definizione di radice quadrata}
% ------------------------------------------------------------

La misura esatta del lato di un quadrato di area \(a\) si indica con il simbolo
\(\sqrt{a}\) e si legge \emph{radice quadrata di \(a\)}.

\begin{definizione}
    La \textbf{radice quadrata} di un numero \(a \geq 0\) è il numero \(x \geq 0\)
    tale che \(x^2 = a\). In simboli:
    \[
        \text{se } a \geq 0 \colon \quad \sqrt{a} = x \iff x^2 = a \quad (x \geq 0)
    \]
    La radice quadrata è l'operazione inversa dell'\textbf{elevamento al quadrato}.
\end{definizione}

\begin{attenzione}
    \textbf{Attenzione alla condizione \(a \geq 0\).}

    La radice quadrata è definita solo per numeri non negativi. Non esiste nessun
    numero reale il cui quadrato sia negativo: qualunque numero moltiplicato per
    sé stesso dà sempre un risultato positivo o zero. Per questo motivo
    \(\sqrt{-4}\) non esiste nell'insieme dei numeri che utilizziamo alle medie.
\end{attenzione}

% ------------------------------------------------------------
\section{I numeri quadrati perfetti}
% ------------------------------------------------------------

\begin{definizione}
    Un numero intero è detto \textbf{quadrato perfetto} se è uguale al quadrato
    di un numero intero. In altre parole, \(n\) è un quadrato perfetto se
    \(\sqrt{n}\) è un numero intero.
\end{definizione}

I primi quadrati perfetti sono:

\begin{center}
    \begin{tabular}{ccc}
        \toprule
        Numero intero \(n\) & Quadrato \(n^2\) & Radice \(\sqrt{n^2}\) \\
        \midrule
        1                   & 1                & 1                     \\
        2                   & 4                & 2                     \\
        3                   & 9                & 3                     \\
        4                   & 16               & 4                     \\
        5                   & 25               & 5                     \\
        6                   & 36               & 6                     \\
        7                   & 49               & 7                     \\
        8                   & 64               & 8                     \\
        9                   & 81               & 9                     \\
        10                  & 100              & 10                    \\
        11                  & 121              & 11                    \\
        12                  & 144              & 12                    \\
        13                  & 169              & 13                    \\
        14                  & 196              & 14                    \\
        15                  & 225              & 15                    \\
        20                  & 400              & 20                    \\
        \bottomrule
    \end{tabular}
\end{center}

Conoscere questi quadrati perfetti a memoria è molto utile per calcolare
rapidamente le radici quadrate esatte.

\begin{esempio}
    Calcola le seguenti radici quadrate esatte e verifica il risultato.

    \medskip

    \begin{tabular}{ll}
        \(\sqrt{16} = 4\)          & perché \(4^2 = 16\) \quad \textit{(corretto)}          \\[4pt]
        \(\sqrt{225} = 15\)        & perché \(15^2 = 225\) \quad \textit{(corretto)}        \\[4pt]
        \(\sqrt{12{,}25} = 3{,}5\) & perché \(3{,}5^2 = 12{,}25\) \quad \textit{(corretto)} \\[4pt]
        \(\sqrt{81} = 9\)          & perché \(9^2 = 81\) \quad \textit{(corretto)}          \\[4pt]
        \(\sqrt{169} = 13\)        & perché \(13^2 = 169\) \quad \textit{(corretto)}        \\[4pt]
        \(\sqrt{64} = 8\)          & perché \(8^2 = 64\) \quad \textit{(corretto)}          \\
    \end{tabular}
\end{esempio}

% ------------------------------------------------------------
\section{Numeri razionali e numeri irrazionali}
% ------------------------------------------------------------

Come abbiamo visto, la radice quadrata di un quadrato perfetto è sempre
un numero intero o un numero decimale finito (come \(3{,}5\)).
Questi numeri appartengono all'insieme dei \textbf{numeri razionali}.

\begin{definizione}
    Un \textbf{numero razionale} è un numero che può essere scritto come
    frazione \(\dfrac{p}{q}\) con \(p\) e \(q\) interi e \(q \neq 0\).
    Nella sua forma decimale, un numero razionale è sempre:
    \begin{itemize}
        \item un decimale finito (es.\ \(3{,}5 = \dfrac{7}{2}\)), oppure
        \item un decimale periodico (es.\ \(0{,}\overline{3} = \dfrac{1}{3}\)).
    \end{itemize}
\end{definizione}

\begin{definizione}
    Un \textbf{numero irrazionale} è un numero che \emph{non} può essere
    scritto come frazione tra due interi. Nella sua forma decimale ha
    infinite cifre dopo la virgola che non si ripetono mai con un periodo
    regolare.

    La radice quadrata di un numero che non è un quadrato perfetto è sempre
    un numero irrazionale.
\end{definizione}

\begin{esempio}
    \(\sqrt{8}\) è un numero irrazionale:
    \[
        \sqrt{8} = 2{,}828\,427\,124\,746\,190\ldots
    \]
    Le cifre decimali continuano all'infinito senza alcuna ripetizione periodica.
    Non possiamo scrivere \(\sqrt{8}\) come frazione esatta; possiamo solo
    \textbf{approssimarlo}.

    Allo stesso modo:
    \[
        \sqrt{11} = 3{,}316\,624\,790\ldots \qquad
        \sqrt{3}  = 1{,}732\,050\,807\ldots
    \]
\end{esempio}

\begin{sapevi}
    \textbf{La crisi dei Pitagorici e \(\sqrt{2}\)}

    I Pitagorici (VI secolo a.C.) credevano che tutti i numeri fossero
    razionali, cioè esprimibili come rapporto di interi. Fu un grande
    shock scoprire che la diagonale di un quadrato di lato 1 — che misura
    esattamente \(\sqrt{2}\) — non è razionale.

    Secondo alcune fonti, il matematico Ippaso di Metaponto fu espulso
    dalla scuola pitagorica (o addirittura annegato!) per aver rivelato
    questa scoperta scomoda. \(\sqrt{2} = 1{,}414\,213\,562\ldots\) rimane
    uno dei numeri irrazionali più famosi della storia della matematica.
\end{sapevi}

% ------------------------------------------------------------
\section{Approssimare la radice quadrata}
% ------------------------------------------------------------

Quando la radice quadrata non è un numero esatto, si effettua
un'\textbf{approssimazione}: ci si accontenta di un certo numero di
cifre decimali, sapendo di commettere un piccolo errore.
Al posto del simbolo \(=\) si usa il simbolo \(\approx\) (che si legge
\emph{circa uguale a}).

\begin{regola}
    \textbf{Regola di arrotondamento}

    Per approssimare un numero decimale a una data cifra:
    \begin{itemize}
        \item se la cifra \emph{successiva} è \(\geq 5\), si \textbf{arrotonda per eccesso}
              (si aumenta di 1 l'ultima cifra che si vuole conservare);
        \item se la cifra \emph{successiva} è \(< 5\), si \textbf{arrotonda per difetto}
              (si elimina semplicemente la parte restante).
    \end{itemize}
\end{regola}

\begin{esempio}
    Approssima \(\sqrt{8}\) e \(\sqrt{11}\) alle cifre indicate.

    \[
        \sqrt{8} = 2{,}828\,427\ldots
    \]
    \begin{itemize}
        \item ai decimi: \(\sqrt{8} \approx 2{,}8\) \quad
              (la cifra successiva è 2, quindi per difetto) \quad \textit{(corretto)}
        \item ai centesimi: \(\sqrt{8} \approx 2{,}83\) \quad
              (la cifra successiva è 8, quindi per eccesso) \quad \textit{(corretto)}
    \end{itemize}

    \[
        \sqrt{11} = 3{,}316\,624\ldots
    \]
    \begin{itemize}
        \item ai decimi: \(\sqrt{11} \approx 3{,}3\) \quad
              (la cifra successiva è 1, quindi per difetto) \quad \textit{(corretto)}
        \item ai centesimi: \(\sqrt{11} \approx 3{,}32\) \quad
              (la cifra successiva è 6, quindi per eccesso) \quad \textit{(corretto)}
    \end{itemize}
\end{esempio}

% ------------------------------------------------------------
\section{Proprietà della radice quadrata}
% ------------------------------------------------------------

\begin{regola}
    \textbf{Proprietà fondamentali della radice quadrata}

    Siano \(a \geq 0\) e \(b \geq 0\) due numeri reali. Valgono le seguenti proprietà:

    \medskip

    \textbf{1. Prodotto sotto radice}
    \[
        \sqrt{a \cdot b} = \sqrt{a} \cdot \sqrt{b}
    \]

    \textbf{2. Quoziente sotto radice} \quad (con \(b > 0\))
    \[
        \sqrt{\dfrac{a}{b}} = \dfrac{\sqrt{a}}{\sqrt{b}}
    \]

    \textbf{3. Radice del quadrato}
    \[
        \sqrt{a^2} = a \quad (a \geq 0)
    \]

    \textbf{4. Quadrato della radice}
    \[
        \left(\sqrt{a}\right)^2 = a \quad (a \geq 0)
    \]
\end{regola}

\begin{esempio}
    Applica le proprietà della radice quadrata.

    \medskip

    \textbf{Proprietà del prodotto:}
    \[
        \sqrt{36} = \sqrt{4 \cdot 9} = \sqrt{4} \cdot \sqrt{9} = 2 \cdot 3 = 6
        \quad \textit{(corretto)}
    \]
    \[
        \sqrt{300} = \sqrt{100 \cdot 3} = \sqrt{100} \cdot \sqrt{3}
        = 10\sqrt{3} \approx 10 \cdot 1{,}732 = 17{,}32
        \quad \textit{(corretto)}
    \]

    \textbf{Proprietà del quoziente:}
    \[
        \sqrt{\frac{25}{4}} = \frac{\sqrt{25}}{\sqrt{4}} = \frac{5}{2} = 2{,}5
        \quad \textit{(corretto)}
    \]

    \textbf{Quadrato della radice:}
    \[
        \left(\sqrt{7}\right)^2 = 7
        \quad \textit{(corretto)}
    \]
\end{esempio}

\begin{attenzione}
    \textbf{Errori comuni da evitare}

    \medskip

    \textbf{Errore 1: somma sotto radice.}

    La radice \emph{non} si distribuisce sulla somma:
    \[
        \sqrt{a + b} \neq \sqrt{a} + \sqrt{b}
    \]
    Esempio: \(\sqrt{9 + 16} = \sqrt{25} = 5\), ma
    \(\sqrt{9} + \sqrt{16} = 3 + 4 = 7 \neq 5\).

    \medskip

    \textbf{Errore 2: radice di un numero negativo.}

    Espressioni come \(\sqrt{-9}\) non hanno significato nell'insieme dei
    numeri reali. Verificare sempre che il numero sotto la radice sia
    non negativo prima di calcolare.

    \medskip

    \textbf{Errore 3: confondere quadrato e radice.}

    \(\sqrt{25} = 5\), non \(5^2 = 25\). Sono operazioni inverse: la radice
    \emph{disfa} il quadrato, il quadrato \emph{disfa} la radice.
\end{attenzione}

% ------------------------------------------------------------
\section{Applicazioni geometriche}
% ------------------------------------------------------------

La radice quadrata compare spesso in problemi di geometria piana,
in particolare quando si conosce l'area di una figura e si vuole
ricavare una lunghezza.

\begin{esempio}
    \textbf{Lato di un quadrato data l'area.}

    Un quadrato ha area \(A = 576\ \text{m}^2\). Qual è la lunghezza del suo lato?
    \[
        \ell = \sqrt{A} = \sqrt{576} = 24\ \text{m}
        \quad \textit{(corretto)}
    \]
    Verifica: \(24^2 = 576\). \quad \textit{(corretto)}
\end{esempio}

\begin{esempio}
    \textbf{Diagonale di un quadrato (Teorema di Pitagora).}

    Per un quadrato di lato \(\ell\), la diagonale \(d\) si calcola con il
    Teorema di Pitagora: \(d^2 = \ell^2 + \ell^2 = 2\ell^2\), quindi:
    \[
        d = \sqrt{2\ell^2} = \ell\sqrt{2}
    \]

    Se il lato è \(\ell = 0{,}6\ \text{m}\):
    \[
        d = 0{,}6 \cdot \sqrt{2} \approx 0{,}6 \cdot 1{,}414 = 0{,}848\ \text{m}
        \approx 0{,}848\ \text{m}
        \quad \textit{(corretto)}
    \]
\end{esempio}

% ------------------------------------------------------------
\section{Esercizi}
% ------------------------------------------------------------

\subsection{Calcolo con la calcolatrice}

\begin{enumerate}[label=\arabic*.]

    \item Calcola le seguenti radici quadrate con la calcolatrice e approssima
          alle cifre indicate.

          \begin{center}
              \begin{tabular}{llll}
                  \toprule
                  Radice            & Ai decimi          & Ai centesimi        & All'unità      \\
                  \midrule
                  \(\sqrt{420}\)    & \(\approx 20{,}5\) & \(\approx 20{,}49\) & \(\approx 20\) \\
                  \(\sqrt{2430}\)   & \(\approx 49{,}3\) & \(\approx 49{,}30\) & \(\approx 49\) \\
                  \(\sqrt{255}\)    & \(\approx 15{,}9\) & \(\approx 15{,}97\) & \(\approx 16\) \\
                  \(\sqrt{270}\)    & \(\approx 16{,}4\) & \(\approx 16{,}43\) & \(\approx 16\) \\
                  \(\sqrt{2{,}45}\) & \(\approx 1{,}6\)  & \(\approx 1{,}56\)  & \(\approx 2\)  \\
                  \(\sqrt{123}\)    & \(\approx 11{,}1\) & \(\approx 11{,}09\) & \(\approx 11\) \\
                  \bottomrule
              \end{tabular}
          \end{center}

    \item Completa la seguente tabella relativa a quadrati:

          \begin{center}
              \begin{tabular}{ccc}
                  \toprule
                  Lato                           & Area                     & Perimetro                      \\
                  \midrule
                  \(8{,}5\ \text{cm}\)           & \(72{,}25\ \text{cm}^2\) & \(34\ \text{cm}\)              \\
                  \(24\ \text{m}\)               & \(576\ \text{m}^2\)      & \(96\ \text{m}\)               \\
                  \(16\ \text{cm}\)              & \(256\ \text{cm}^2\)     & \(64\ \text{cm}\)              \\
                  \(\approx 22{,}98\ \text{dm}\) & \(528\ \text{dm}^2\)     & \(\approx 91{,}92\ \text{dm}\) \\
                  \bottomrule
              \end{tabular}
          \end{center}

    \item Calcola la misura del lato di un quadrato nel caso in cui la sua
          diagonale misuri \(15\ \text{cm}\). Approssima tutti i risultati al
          decimo.

    \item Calcola la misura della diagonale di un quadrato nel caso in cui il
          lato misuri \(0{,}6\ \text{m}\). Approssima tutti i risultati al
          millesimo.

    \item L'area di un triangolo rettangolo isoscele vale \(4{,}5\ \text{cm}^2\).
          Trova la misura dei due lati congruenti.

    \item L'area della superficie totale di un cubo è \(324\ \text{cm}^2\).
          Calcola il suo volume.

\end{enumerate}

% ------------------------------------------------------------
\section{Riepilogo}
% ------------------------------------------------------------

\begin{itemize}
    \item La \textbf{radice quadrata} di \(a \geq 0\) è il numero \(x \geq 0\)
          tale che \(x^2 = a\): \(\sqrt{a} = x \iff x^2 = a\).
    \item La radice quadrata è l'\textbf{operazione inversa} dell'elevamento
          al quadrato.
    \item I \textbf{quadrati perfetti} sono i numeri interi il cui la radice
          quadrata è anch'essa un numero intero (1, 4, 9, 16, 25, \ldots).
    \item Un numero razionale si scrive come frazione o decimale finito/periodico.
          Un \textbf{numero irrazionale} ha infinite cifre decimali non periodiche;
          \(\sqrt{2}\), \(\sqrt{3}\), \(\sqrt{8}\) sono esempi di numeri irrazionali.
    \item Quando la radice non è esatta si \textbf{approssima} usando \(\approx\),
          rispettando la regola di arrotondamento (\(\geq 5\) per eccesso,
          \(< 5\) per difetto).
    \item \textbf{Proprietà:} \(\sqrt{a \cdot b} = \sqrt{a}\cdot\sqrt{b}\);
          \(\sqrt{a/b} = \sqrt{a}/\sqrt{b}\); \((\sqrt{a})^2 = a\).
    \item \textbf{Non vale} \(\sqrt{a+b} = \sqrt{a}+\sqrt{b}\).
    \item La radice quadrata non esiste (nei reali) per numeri negativi.
\end{itemize}

% ------------------------------------------------------------
\section{Soluzioni degli esercizi}
% ------------------------------------------------------------

\subsection*{Esercizio 3}

La diagonale di un quadrato di lato \(\ell\) vale \(d = \ell\sqrt{2}\),
quindi \(\ell = d / \sqrt{2}\).

\[
    \ell = \frac{15}{\sqrt{2}} = \frac{15}{1{,}414\ldots} \approx 10{,}6\ \text{cm}
\]

\subsection*{Esercizio 4}

\[
    d = 0{,}6 \cdot \sqrt{2} \approx 0{,}6 \cdot 1{,}414 = 0{,}848\ \text{m}
\]

\subsection*{Esercizio 5}

Un triangolo rettangolo isoscele ha i due cateti uguali (\(\ell\)).
L'area è \(A = \ell^2 / 2\), quindi:
\[
    \ell^2 = 2A = 2 \cdot 4{,}5 = 9 \implies \ell = \sqrt{9} = 3\ \text{cm}
\]

\subsection*{Esercizio 6}

La superficie totale di un cubo di spigolo \(a\) è \(S = 6a^2\).
\[
    a^2 = \frac{S}{6} = \frac{324}{6} = 54\ \text{cm}^2
    \implies a = \sqrt{54} \approx 7{,}35\ \text{cm}
\]
Il volume è:
\[
    V = a^3 \approx (7{,}35)^3 \approx 396{,}9\ \text{cm}^3
\]